\documentclass[11pt]{article}
\usepackage[a4paper,margin=28mm]{geometry}
\usepackage{amsmath,amssymb,amsthm,mathtools}
\usepackage{mathrsfs}
\usepackage{booktabs,longtable,array,tabularx}
\usepackage{xcolor}
\usepackage{microtype}
\usepackage[colorlinks=true,linkcolor=blue!60!black,citecolor=blue!60!black,urlcolor=blue!65!black]{hyperref}
\usepackage[nameinlink,noabbrev]{cleveref}
\setlength{\emergencystretch}{2em}

\newtheorem{theorem}{Theorem}[section]
\newtheorem{proposition}[theorem]{Proposition}
\newtheorem{lemma}[theorem]{Lemma}
\newtheorem{corollary}[theorem]{Corollary}
\newtheorem{hypothesis}[theorem]{Hypothesis}
\theoremstyle{definition}
\newtheorem{definition}[theorem]{Definition}
\newtheorem{remark}[theorem]{Remark}
\newcommand{\ThetaN}{\Theta(N)}
\newcommand{\cA}{\mathcal A}
\newcommand{\cP}{\mathcal P}
\newcommand{\cG}{\mathcal G}
\newcommand{\cT}{\mathcal T}
\newcommand{\Ups}{\Upsilon}
\newcommand{\eps}{\varepsilon}
\newcommand{\R}{\mathbb R}
\newcommand{\Q}{\mathbb Q}
\newcommand{\openstatus}{\textcolor{red!75!black}{OPEN}}
\newcommand{\closedstatus}{\textcolor{green!45!black}{CLOSED}}
\newcommand{\conditionalstatus}{\textcolor{orange!80!black}{CONDITIONAL}}
\newcommand{\cert}[1]{\texttt{#1}}

\title{An Audited Conditional Wu--BJS Assembly at the\\
Double-Exponential Threshold $c=23.36$}
\author{Revised proof record}
\date{2 September 2026}

\begin{document}
\maketitle

\begin{abstract}
This is a proof-audit revision of the manuscript claiming an explicit Chen
theorem at \(c=23.36\).  The revision corrects the local-density convention,
the finite Wu bad-set argument, the tagged linear-sieve main term, the
five-box endpoints, the Johnston notation, and the numerical ledger.
It also records the exact hypotheses needed for the \(u_0=10^7\) product
condition, the \(X\geq N^{0.398}\) rectangle extension, and the
\(U_{13},U_{14}\) complementary-prime switch.  Those three interfaces are
not consequences of the published BJS statements as quoted here, and no
certificate proving them was supplied with the input files.  Consequently
this document does not claim an unconditional Chen theorem.  Its strongest
strict conclusion is a conditional assembly theorem: if the six explicitly
listed interfaces are subsequently proved with the stated quantifiers, then
the old endpoint ledger would imply the advertised positive lower bound.
\end{abstract}

\tableofcontents

\section{Scope and conclusion}
\label{sec:scope}

The sole audit standard for this revision is the supplied \cert{reply.md}.
That file is treated as an audit report, not as a mathematical source and not
as an instruction embedded in a theorem.  The external mathematical sources
used below are the published BJS paper \cite{BJS} and Wu's fifteen-row
decomposition \cite{WuI}.  Every statement not proved in this file is labelled
as quoted, conditional, or open.

Put
\[
 c=\log\log N,\qquad L=\log N=e^c,\qquad
 N_0=\exp(\exp(23.36)).
\]
For an even \(N\), write
\[
 D_{1,2}(N)=
 \#\{p\leq N:p\ {\rm prime},\ N-p>0,\ \Omega(N-p)\leq2\}.
\]

\begin{theorem}[Strongest conclusion of this revision]
\label{thm:conditional-conclusion}
Assume the six hypotheses
\(\mathrm{H}_{\rm prod}\), \(\mathrm{H}_{\rm rect}\),
\(\mathrm{H}_{\rm sw}\), \(\mathrm{H}_{\rm exc}\),
\(\mathrm{H}_{\rm led}\), and \(\mathrm{H}_{\rm half}\) stated in
\cref{sec:critical-hypotheses}.  Assume also the quoted BJS and Wu results
in \cref{sec:external-inputs} in exactly the stated versions.  Then every
even \(N\geq N_0\) satisfies
\[
 D_{1,2}(N)\geq 0.1731658\,\Theta(N)>0.17\,\Theta(N).
\]
\end{theorem}

\begin{proof}
The corrected Wu inequality is \cref{lem:wu-corrected}.  The ordinary rows
are bounded only after the tagged application of the BJS sieve in
\cref{lem:tagged}; the five-box rows require \(\mathrm{H}_{\rm rect}\);
and the two negative switched rows require \(\mathrm{H}_{\rm sw}\).
The exceptional and terminal charges are covered by
\(\mathrm{H}_{\rm exc}\) and the terminal part of \(\mathrm{H}_{\rm led}\).
The single ledger inequality in \(\mathrm{H}_{\rm led}\) gives the endpoint
margin, and \(\mathrm{H}_{\rm half}\) extends that inequality over the
whole half-line.  No other step is used.
\end{proof}

\begin{proposition}[Unconditional status of the file]
\label{prop:status}
Without the six hypotheses in the preceding theorem, this file proves only
the finite algebraic identities, the corrected elementary terminal bounds,
the exact interval-arithmetic assertions, and the necessity of the listed
interfaces.  It does not prove \(D_{1,2}(N)>0\) for any whole half-line
\(c\geq23.36\).  Thus the final classification is:
\[
\boxed{\begin{gathered}
\text{\rm Class 3: a core analysis interface remains unproved;}\\
\text{\rm only a conditional reduction theorem is retained.}
\end{gathered}}
\]
\end{proposition}

\section{Definitions and parameter conventions}
\label{sec:definitions}

\subsection{Local products and normalization}
\label{sec:products}

Define
\[
 C_N=
 \prod_{\substack{q\mid N\\q>2}}\frac{q-1}{q-2}
 \prod_{q>2}\left(1-\frac1{(q-1)^2}\right),
 \qquad
 \Theta(N)=\frac{C_NN}{L^2}.
\]
The standard lower bound \(C_N>0.6601617>0.66\) is used only as a
normalization bound.  For \(u\ge2\), define the set and the product
\[
 \mathbb P_N=\{p:p\ {\rm prime},\ p\nmid N\},\qquad
 \mathbb P_N(u)=\{p\in\mathbb P_N:p<u\},\qquad
 Q_N(u)=\prod_{p\in\mathbb P_N(u)}p.
\]
The full product, used only after an explicit coprimality restriction, is
\[
 P^{\rm all}(u)=\prod_{p<u}p.
\]
Thus \(Q_N(u)\) is not an undefined synonym for a full primorial.  The
worst fixed envelope is
\[
 Q_{\max}(u)=P^{\rm all}(u),\qquad Q_N(u)\mid Q_{\max}(u),
 \qquad \log Q_N(u)\leq\log Q_{\max}(u).
\]
For a source tag \(\alpha\), its local density is denoted \(g_\alpha\).
The convention used in all conditional sieve statements is
\[
 g_\alpha(p)=
 \begin{cases}
 0,&p\mid \nu_\alpha,\\[2pt]
 1/(p-1),&p\nmid \nu_\alpha,
 \end{cases}
\tag{2.1}
\]
for odd primes \(p\), where \(\nu_\alpha\) is the fixed part of the tag.
The prime \(2\) is absent from the odd source sequences; it is never silently
inserted into an odd-prime product.

\subsection{Fixed parameters}
\label{sec:parameters}

\[
\begin{array}{c|cccccccc}
 &\kappa_1&\kappa_2&\rho&\sigma_2&\sigma_1&u_0&\delta&\eps_0\\
\hline
\text{value}&1/12&29/250&1/4&141/500&41/125&10^7&29/25&1/10000
\end{array}
\tag{2.2}
\]
The support parameter in \(Z\) and the BJS product tolerance are different
numbers:
\[
 \eps=\frac1{2000},\qquad
 \eta=\frac{1452}{10^{10}},\qquad
 C_1(\eta)\le106,\qquad C_2(\eta)\le107.
\tag{2.3}
\]
Here \(\eps\) occurs only in the support definition (2.5), whereas
\(\eta\) is the tolerance in the linear-sieve product hypothesis.  The
bounds for \(C_1(\eta),C_2(\eta)\) follow conservatively from the BJS table
at tolerance \(10^{-5}\), since these constants are nonincreasing as the
reciprocal tolerance increases.  None of \(\eps,\eta,\eps_0\) may be
interchanged.
The main scales are
\[
 y=N^{1/3},\qquad z=N^{1/8},\qquad z_*=N^{\kappa_1},
\qquad X_0=N^{1/8},
\tag{2.4}
\]
\[
 Z=N^{1/(1+\eps/2)},\qquad
 z_0=Z^{1/(\log\log Z)^3},\qquad
 D_\circ=z_0^3Z^{1969/12000}.
\tag{2.5}
\]
For later use define distinct names for the two quantities previously both
called \(H\):
\[
 H_{\rm en}:=\text{the coefficient-energy envelope},\qquad
 H_{\rm lev}(c):=\frac{20\log(L/12)}{L/12}.
\tag{2.6}
\]
The only numerical energy claim that may be used is an explicitly certified
bound \(H_{\rm en}\le0.392000001\).  The level quantity
\(H_{\rm lev}(23.36)\) is of order \(10^{-7}\), and is not that energy
constant.

\subsection{The source domains}
\label{sec:domains}

The two difficult four-prime domains are
\[
\begin{aligned}
 \mathscr T_{13}&=\{(p_1,p_2,p_3,p_4):
 N^{\kappa_1}\le p_1<p_2<p_3<p_4<N^{\kappa_2}\},\\
 \mathscr T_{14}&=\{(p_1,p_2,p_3,p_4):
 N^{\kappa_1}\le p_1<p_2<p_3<N^{\kappa_2}\le p_4
 <N^{1/2-2\kappa_1}/p_3\}.
\end{aligned}
\tag{2.7}
\]
All source terms also carry \((p_1p_2p_3p_4,N)=1\).  A source element
in a box is always defined with the positivity condition \(pn<N\); this
condition is not inferred from \(p<(1+\eps_0)\omega\) and \(n<N/\omega\).

\section{External inputs and their exact scope}
\label{sec:external-inputs}

\subsection{The BJS linear sieve}
\label{sec:bjs-linear}

The following is the form of BJS Theorem 6 that is used here.  Let
\(A=\{a(n)\}_{n\ge1}\) be a finite nonnegative sequence, let
\(\mathbb P\) be a set of primes, and set
\[
 P(z)=\prod_{\substack{p\in\mathbb P\\p<z}}p,\qquad
 S(A,\mathbb P,z)=\sum_{(n,P(z))=1}a(n).
\]
For every tag \(n\), let \(g_n\) be multiplicative with
\(0\le g_n(p)<1\).  If, for some \(0<\xi\le1/74\),
\[
 A_d:=\sum_{d\mid n}a(n)
 =\sum_n a(n)g_n(d)+r(d),
\tag{3.1}
\]
and a finite set \(\mathscr Q\subseteq\mathbb P\) has product
\(Q=\prod_{q\in\mathscr Q}q\), BJS assumes, for every tag
and every \(1<u<z\),
\[
 \prod_{\substack{p\in\mathbb P\setminus\mathscr Q\\u\le p<z}}
 (1-g_n(p))^{-1}
 <(1+\xi)\frac{\log z}{\log u}.
\tag{3.2}
\]
With
\[
 X_A=\sum_n a(n)\prod_{p\mid P(z)}(1-g_n(p)),\qquad
 R=\sum_{\substack{d\mid P(z)\\d<QD}}|r(d)|,\qquad
 s=\frac{\log D}{\log z},
\tag{3.3}
\]
BJS gives
\[
 S(A,\mathbb P,z)
 <\bigl(F(s)+\xi C_1(\xi)e^2h(s)\bigr)X_A+R
\quad(D\ge z),
\tag{3.4}
\]
and
\[
 S(A,\mathbb P,z)
 >\bigl(f(s)-\xi C_2(\xi)e^2h(s)\bigr)X_A-R
\quad(D\ge z^2),
\tag{3.5}
\]
where
\[
 h(s)=
 \begin{cases}
 e^{-2},&1\le s\le2,\\
 e^{-s},&2\le s\le3,\\
 3s^{-1}e^{-s},&s\ge3.
 \end{cases}
\tag{3.6}
\]
The functions \(F,f\) are the BJS delay functions, with \(f(s)=0\) for
\(s\le2\).  The tag-dependent quantity \(X_A\) in (3.3) is not replaced by
\(|A|V(z)\) unless all tags have the same \(g_n\).

\begin{hypothesis}[Finite product certificate at \(u_0=10^7\)]
\label{hyp:prod}
\(\mathrm{H}_{\rm prod}\) is the following quantified assertion.  For every
even \(N\ge N_0\), every tag \(\alpha\) occurring in the source partitions,
and every pair \(1<u<z\) for which (3.4) or (3.5) is invoked,
\[
 \prod_{\substack{p\in\mathbb P_N\setminus\mathbb P_N(u_0)\\u\le p<z}}
 \bigl(1-g_\alpha(p)\bigr)^{-1}
 <(1+\eta)\frac{\log z}{\log u},
 \qquad \eta=\frac{1452}{10^{10}}.
\tag{3.7}
\]
In applying (3.2)--(3.5), one must take \(\xi=\eta\); the unrelated
support parameter \(\eps=1/2000\) is not substituted for \(\xi\).
The same certificate must give an outward interval for
\(\log Q_N(u_0)\) and a fixed worst-case interval obtained from
\[
 Q_N(u_0)\mid P^{\rm all}(u_0).
\tag{3.8}
\]
The proof record for (3.7) must contain: (i) the deterministic list of all
primes \(p\le u_0\), with a hash; (ii) the exact finite product or
log-product accumulated with directed rational rounding; (iii) an explicit
Abel summation identity for \(p>u_0\); (iv) an analytic tail bound with its
sign and constant; and (v) a quantifier check over every allowed \(u,z,N\)
and every tag.  For example, if
\(b(t)=\log(1+1/(t-2))\), the tail must be bounded from an explicitly
stated \(\vartheta(t)\) or \(\pi(t)\) envelope by
\[
 \sum_{u_0<p<z}b(p)
 =A(z)b(z)-A(u_0)b(u_0)
 -\int_{u_0}^{z}A(t)b'(t)\,dt,
\tag{3.9}
\]
with \(A(t)\) and the endpoint convention defined in the certificate.
\end{hypothesis}

\begin{remark}
The two decimal tail values \(1.2558668307\cdot10^{-7}\) and
\(1.0727189017\cdot10^{-7}\) in the former manuscript are not a proof of
(3.7): they do not specify the finite product, prime segmentation, Abel
identity, directed rounding, or source hash.  Thus the \(u_0=10^7\) input is
\openstatus, even though its exact target is now stated.
\end{remark}

\subsection{The published rectangle lemma}
\label{sec:bjs-rectangle}

BJS Lemma 32 is quoted here because it fixes the exact range that the
revision is allowed to use.  Let \(y=N^{1/3}\), \(z=N^{1/8}\).  It assumes
\[
 \frac Ny<X\le\frac Nz,\qquad XY<2N,\qquad Y>Z,\qquad Y>z,
\tag{3.7}
\]
as well as \(k_2<K_\delta(x_2(Y))\), \(N>(X_3)^8\), and
\(\log\log x_2(X_3)\ge10.4\).  With
\[
 D^*=\frac{\sqrt{XY}}{\log^{10}Y},
\tag{3.8}
\]
it bounds the prime--integer rectangle discrepancy by
\[
 \sum_{\substack{d<D^*\\d\mid P(y)}}\max_{(a,d)=1}
 \left|
 \sum_{n<X}\sum_{\substack{Z\le p<Y\\np\equiv N\pmod d}}a(n)
 -\frac1{\varphi(d)}
 \sum_n\sum_{\substack{Z\le p<Y\\(np,d)=1}}a(n)
 \right|
 \le \frac{m(X_3)XY}{\log^3Y},
\tag{3.9}
\]
for \(|a(n)|\le1\), where
\[
\begin{aligned}
 m(X_3)={}&39v_0(X_3)
 +\frac{108\log^{16}X_3}{X_3\log\log X_3}
 +\frac{26\log^5X_3}{\log^{10}(x_2(X_3))}\\
 &+88\log^5X_3\left(X_3^{-8/3}+X_3^{-1/2}\right)
 +\frac{106}{\log^6X_3}.
\end{aligned}
\tag{3.10}
\]
Since \(N/y=N^{2/3}\), (3.7) does not cover the boxes below
\(N^{2/3}\).  The extension to \(X\ge N^{0.398}\) is therefore a new
interface, not a quotation.

\subsection{Wu's signed source identity}
\label{sec:wu-input}

We quote Wu's fifteen-row identity in the following sign convention:
\[
\begin{aligned}
4D_{1,2}(N)\ge{}&
4S(\cA;\mathbb P_N,z_*)-\Ups_1-\Ups_2-\Ups_3
 +\Ups_4+\Ups_5+\Ups_6\\
&-2\Ups_7-2\Ups_8-\Ups_9-\Ups_{10}
 +\Ups_{11}+\Ups_{12}-\Ups_{13}-\Ups_{14}+\Ups_{15}.
\end{aligned}
\tag{3.11}
\]
This is an external combinatorial input.  The finite extension from its good
source to all sources is proved below; the four positive rows
\(\Ups_6,\Ups_{11},\Ups_{12},\Ups_{15}\) are deliberately removed before
that extension.

\section{The corrected finite Wu reduction}
\label{sec:wu}

\subsection{Good and bad source elements}
\label{sec:bad}

For an odd prime \(p\), put \(a=N-p\).  Let \(\cG_N\) be the submultiset
of \(a\)'s that are square-free, coprime to \(N\), and free of primes below
\(z_*\).  The source with \(p=2\) is kept in a separate one-element charge.
On \(\cG_N\), Wu's Buchstab identities are finite exact identities.

\begin{lemma}[Corrected finite Wu inequality]
\label{lem:wu-corrected}
For every even \(N\ge N_0\),
\[
\begin{aligned}
4D_{1,2}(N)\ge{}&
4S(\cA;\mathbb P_N,z_*)-\Ups_1-\Ups_2-\Ups_3
 +\Ups_4+\Ups_5\\
&-2\Ups_7-2\Ups_8-\Ups_9-\Ups_{10}
 -\Ups_{13}-\Ups_{14}-E_{\rm bad}(N),
\end{aligned}
\tag{4.1}
\]
where
\[
 E_{\rm bad}(N)=
 \left(4+\left(\frac L{\log2}\right)^2\right)
 \left(2N^{1-\kappa_1}+\pi(\sqrt N)+\frac L{\log2}+1\right).
\tag{4.2}
\]
\end{lemma}

\begin{proof}
Apply (3.11) on \(\cG_N\) and delete the positive terms
\(\Ups_6,\Ups_{11},\Ups_{12},\Ups_{15}\).  Deleting positive terms lowers
the right-hand side and is therefore valid before any bad-set compensation.

If \(q^2\mid a\) with \(q\ge N^{\kappa_1}\), then \(q\le\sqrt N\), and the
number of primes \(p\le N\) in the congruence \(p\equiv N\pmod {q^2}\) is at
most \(N/q^2+1\).  Hence the total number of square-divisibility
incidences is at most
\[
 \sum_{N^{\kappa_1}\le q\le\sqrt N}\left(\frac N{q^2}+1\right)
 \le 2N^{1-\kappa_1}+\pi(\sqrt N).
\tag{4.3}
\]
The \(+1\) term is retained; it is not hidden in the first summand.

If \(q\mid N\) and \(q\mid N-p\), then \(q\mid p\), so \(p=q\).  There are
at most \(\omega(N)\le L/\log2\) such source elements.  In the retained
positive rows \(\Ups_4,\Ups_5\), a fixed bad \(a\) can be selected only by
choosing at most two prime divisors of \(a\); its total multiplicity is at
most \(2\binom{\omega(a)}2\).  The base row contributes at most \(4\).
Since \(\omega(a)\le L/\log2\), the retained positive multiplicity is
at most \(4+(L/\log2)^2\).  The \(p=2\) source contributes the final \(+1\).
Multiplying this coefficient by (4.3) and the \(q\mid N\) count gives
(4.2).  No positive high-multiplicity row is compensated, because it was
deleted at the first step.
\end{proof}

\begin{remark}
The earlier terminal \(2N^{1-\kappa_1}\) omitted both the congruence
floor term and the positive rows of multiplicity three and four.  Formula
(4.2) is the corrected elementary bound.  Its normalized size is not used
to close the main theorem until it is inserted into the single ledger.
\end{remark}

\section{Tagged sieve and the five boxes}
\label{sec:boxes}

\subsection{The tagged main term}
\label{sec:tagged}

\begin{lemma}[Tagged finite linear sieve]
\label{lem:tagged}
Let \(\{(b_\alpha,w_\alpha)\}_{\alpha\in\mathcal I}\) be a finite multiset
with \(w_\alpha\ge0\).  Let \(g_\alpha\) be the tag-dependent densities and
define
\[
 A_d=\sum_{\alpha:d\mid b_\alpha}w_\alpha,\qquad
 X_A=\sum_\alpha w_\alpha
 \prod_{p\mid P(z)}(1-g_\alpha(p)),\qquad
 r(d)=A_d-\sum_\alpha w_\alpha g_\alpha(d).
\tag{5.1}
\]
If every tag satisfies the same product condition (3.2), then the BJS
upper and lower bounds apply with this \(X_A\) and with
\[
 R=\sum_{\substack{d\mid P(z)\\d<QD}}|r(d)|.
\tag{5.2}
\]
\end{lemma}

\begin{proof}
Represent the multiset by the nonnegative sequence
\(a(b_\alpha)\) with repeated tags retained.  Then (5.1) is exactly (3.1)
and (5.2) is exactly (3.3).  The BJS Rosser coefficients depend on \(D,z,Q\),
not on the name of the tag, so the finite sum is
\[
 \sum_d\lambda^\pm(d)A_d
 =\sum_\alpha w_\alpha\sum_d\lambda^\pm(d)g_\alpha(d)
 +\sum_d\lambda^\pm(d)r(d).
\tag{5.3}
\]
Applying the quoted BJS inequality to the tag sequence gives the stated
formula.  No aggregated function
\(\sum_\alpha w_\alpha g_\alpha\) is declared multiplicative, and no
common \(V(z)\) is introduced unless all \(g_\alpha\) agree.
\end{proof}

\subsection{Triple reduction and positivity}
\label{sec:triple}

The exponent gaps used for the four upper rows are
\[
4\sigma_1-1=0.312,\quad 4\sigma_2-1=0.128,\quad
\kappa_1+3\sigma_1-1=0.067\overline3,\quad
\kappa_2+\sigma_2+2\sigma_1-1=0.054.
\tag{5.4}
\]
Thus two surviving cofactor primes would force one of the four exponents in
(5.4) above \(1\).  The only remaining source is a three-prime tuple plus
explicit boundary charges.  A conservative symbolic charge is
\[
 E_{\rm tri}(N)=
 32\left(2N^{1-\kappa_1}+\pi(\sqrt N)+1\right)(1+L)
 +16\left(\frac L{\log2}\right)^3+6N^{2359/3000}.
\tag{5.5}
\]
This expression is a charge, not an assertion that it has already been
absorbed into a numerical reserve.

\subsection{Half-open boxes and the negative-source repair}
\label{sec:box-definition}

For an exponent interval \([a_i,b_i)\), define
\[
 M_i(c)=\left\lceil
 \frac{(b_i-a_i)e^c}{\log(1+\eps_0)}
 \right\rceil,\qquad
 \omega_{i,j}=N^{a_i}(1+\eps_0)^j,\qquad
 I_{i,j}=[\omega_{i,j},(1+\eps_0)\omega_{i,j}).
\tag{5.6}
\]
The five families are
\[
\begin{array}{c|c|c}
\text{family}&\text{selected coordinate}&[a_i,b_i)\\
\hline
7&p_1&[\sigma_1,1/3)\\
8&p_1&[\sigma_2,1/3)\\
9h&p_1&[1/8,\sigma_1)\\
9\ell&p_2&[\sigma_1,11/24)\\
10&p_3&[\sigma_1,1-\kappa_2-\sigma_2)
\end{array}
\tag{5.7}
\]
For a selected prime \(p\in I_{i,j}\), put
\[
 X_{i,j}=N/\omega_{i,j},\qquad
 Z_{i,j}=\omega_{i,j},\qquad
 Y_{i,j}=(1+\eps_0)\omega_{i,j}.
\tag{5.8}
\]
The uncut rectangle and the positive source are, respectively,
\[
\begin{aligned}
 \mathcal B^0_{i,j}
 &=
 \{N-pn:Z_{i,j}\le p<Y_{i,j},\ n<X_{i,j},\
 a_{i,j}(n)=1\},\\
 \mathcal B^+_{i,j}
 &=
 \{N-pn:Z_{i,j}\le p<Y_{i,j},\ n<X_{i,j},\
 pn<N,\ a_{i,j}(n)=1\}.
\end{aligned}
\tag{5.9}
\]
The strip
\[
 \mathcal S_{i,j}=\{N-pn:Z_{i,j}\le p<Y_{i,j},\
 n<X_{i,j},\ pn\ge N,\ a_{i,j}(n)=1\}
\tag{5.10}
\]
satisfies \(\mathcal B^0_{i,j}=\mathcal B^+_{i,j}\mathbin{\dot\cup}
\mathcal S_{i,j}\).  All sieve sequences in this revision use
\(\mathcal B^+_{i,j}\); the strip is charged separately.  This removes the
possibility of a negative source value.

\begin{lemma}[Half-open box map]
\label{lem:boxes}
Each prime \(p\) in a selected interval belongs to exactly one box in
(5.6).  The five families in (5.7) have, at \(c_0=23.36\), the following
integer counts:
\[
\begin{array}{c|r|r}
\text{family}&M_i&\text{certified quotient interval}\\
\hline
7&744\,971\,295\,075&
[744971295074.63,\ 744971295074.65]\\
8&7\,170\,348\,715\,094&
[7170348715093.37,\ 7170348715093.40]\\
9h&28\,355\,469\,918\,779&
[28355469918778.36,\ 28355469918778.39]\\
9\ell&18\,205\,236\,023\,387&
[18205236023386.43,\ 18205236023387.45]\\
10&38\,272\,900\,284\,460&
[38272900284459.46,\ 38272900284459.49]
\end{array}
\tag{5.11}
\]
Consequently
\[
 M_{\rm actual}=92\,748\,926\,236\,795,\qquad
 M_{\rm out}=M_{\rm actual}+5=92\,748\,926\,236\,800.
\tag{5.12}
\]
\end{lemma}

\begin{proof}
The logarithm in (5.6) is strictly increasing, so the floor of
\(\log(p/N^{a_i})/\log(1+\eps_0)\) is unique.  The five intervals are
half-open and therefore do not duplicate their common endpoints.  The
quotient intervals in (5.11) are outward rational intervals; each lies
strictly between two consecutive integers, so its ceiling is the displayed
\(M_i\).  Summing the five integers gives (5.12).  The five additional
units are an outward remainder count only; they are not source mass.
\end{proof}

\subsection{The accepted global rectangle interface}
\label{sec:rectangle-interface}

For the local tag \(a_{i,j}(n)\), use
\[
 g_{i,j}(q)=
 \begin{cases}
 0,&q\mid pn,\\
 1/(q-1),&q\nmid pn.
 \end{cases}
\tag{5.13}
\]
The congruence discrepancy is
\[
\mathcal R_{i,j}(d)=
 \sum_{n<X_{i,j}}a_{i,j}(n)
 \left(
 \sum_{\substack{Z_{i,j}\le p<Y_{i,j}\\np\equiv N\pmod d}}1
 -\frac1{\varphi(d)}
 \sum_{\substack{Z_{i,j}\le p<Y_{i,j}\\(np,d)=1}}1
 \right).
\tag{5.14}
\]
The required five-box quantity is the single global sum
\[
 \mathcal R_{\rm five}=
 \sum_{i,j}\sum_{\substack{d<Q_N(u_0)D^*_{i,j}\\d\mid P^{\rm all}(y)}}
 \lambda_{i,j}(d)\mathcal R_{i,j}(d),
\tag{5.15}
\]
including the strip charges (5.10), endpoint rows, order diagonals, and any
moduli sharing a source prime.

\begin{hypothesis}[Fixed-cutoff rectangle extension]
\label{hyp:rect}
\(\mathrm{H}_{\rm rect}\) is the assertion that for every even \(N\ge N_0\),
every five-box tag in (5.7), and every permitted source coefficient,
(5.15) has the BJS five-term envelope
\[
 E_{\rm rect}=E_1+E_2+E_3+E_4+E_5,
\tag{5.16}
\]
where
\[
\begin{aligned}
 E_1&=\frac{39w_*XY}{\log^3Y},&
 E_2&=\frac{108X\log^{13}Y}{\log\log Y},\\
 E_3&=\frac{26XY\log^2Y}{\log^{10}x_2(X_0)},&
 E_4&=88XY\log^2Y(X^{-1/2}+Y^{-1/2}),\\
 E_5&=\frac{106XY}{\log^9Y}.
\end{aligned}
\tag{5.17}
\]
Here
\[
 x_2(t)=t/\log^{15}t,\quad
 D_0=\log^{10}x_2(X_0),\quad X_3=X_0/2,
\tag{5.18}
\]
\[
 \rho_0=\frac{\eps_0}{1+\eps_0}
 \frac{(\log X_0+\log(1+\eps_0))^5}{\log^6X_0},\quad
 w_*=(1+\rho_0/2)v'_*(X_0),
\tag{5.19}
\]
and \(v'_*(X_0)\) is the BJS envelope with
\(\beta=1-\nu_{\max}\),
\[
\nu_{\max}=
\min\left\{
\frac{100}{\sqrt{K_\delta(u_{\max})}\log^2K_\delta(u_{\max})},
\frac1{2R_1\log Q_1(u_{\max})}
\right\},
\quad R_1=2.0452.
\tag{5.20}
\]
The assertion includes \(X\ge N^{0.398}\), \(Y\ge N^{1/8}\),
\(XY\le(1+\eps_0)N\), fixed \(D_0\), the ordinary and exceptional
deletion versions, and a proof that the global bound (5.15), rather than
\(M_{\rm out}\max E_{\rm rect}\), is the quantity being charged.
\end{hypothesis}

\begin{remark}
\(\mathrm{H}_{\rm rect}\) is not proved here.  The published BJS condition
(3.7) has \(X>N^{2/3}\), while the lower boxes have only
\(X\ge N^{0.398}\).  The missing proof must list every use of
\(X>N^{2/3}\), \(D_0<D^*\), \(D^*\le Y^4\), \(Y>N^{1/8}\), and the
exceptional deletion.  A single endpoint calculation cannot replace this
row-by-row implication.  The five terms in (5.16) are written with explicit
plus signs; they are not a product of five fractions.
\end{remark}

\section{The corrected complementary-prime switch}
\label{sec:switch}

\subsection{The local-coprimality switch}
\label{sec:rough-switch}

For a difficult tuple put
\[
 \varpi=N-p_1p_2p_3p_4\rho.
\]
The correct local switch is
\[
\begin{aligned}
 &\mathbf1_{\{\varpi\ {\rm prime}\}}
 \mathbf1_{\{(\rho,Q_N(p_2))=1\}}\\
 &\quad\le
 \mathbf1_{\{(\rho,N)=1\}}
 \mathbf1_{\{(\rho,P^{\rm all}(z_*))=1\}}
 W^+_{z_*}(\varpi)
 +\mathbf1_{\{\varpi\mid N\}},
\end{aligned}
\tag{6.1}
\]
for \(\varpi\ge z_*\), with the range \(\varpi<z_*\) separately charged.
Indeed, if the left side is one but \((\rho,N)>1\), take a prime
\(q\mid(\rho,N)\).  Since \(q\mid N-p_1p_2p_3p_4\rho=\varpi\) and
\(\varpi\) is prime, \(\varpi=q\mid N\).  Otherwise
\((\rho,N)=1\) and roughness relative to \(Q_N(p_2)\), with \(p_2\ge z_*\),
implies roughness relative to the full product \(P^{\rm all}(z_*)\).
The upper Rosser weight then gives (6.1).

\begin{lemma}[Explicit count for the \(\varpi\mid N\) terminal]
\label{lem:varpi-N}
The exceptional term in (6.1) has the bound
\[
 T_{\varpi\mid N}\le
 \frac L{\log2}
 \left(N^{4\kappa_2}+N^{3\kappa_2+1/4}\right).
\tag{6.2}
\]
Consequently
\[
 \frac{T_{\varpi\mid N}}{\Theta(N)}
 \le \frac{L^3}{0.66\log2}
 \left(N^{4\kappa_2-1}+N^{3\kappa_2-3/4}\right),
\tag{6.3}
\]
which is decreasing as a function of \(c\) on \(c\ge23.36\).
\end{lemma}

\begin{proof}
There are at most \(L/\log2\) prime divisors \(\varpi\) of \(N\).
For \(\mathscr T_{13}\), the four source primes have fewer than
\(N^{4\kappa_2}\) ordered choices.  For \(\mathscr T_{14}\),
\(p_1,p_2,p_3<N^{\kappa_2}\) and
\[
 p_4<N^{1/2-2\kappa_1}/p_3\le N^{1/2-3\kappa_1}=N^{1/4},
\]
so there are fewer than \(N^{3\kappa_2+1/4}\) choices.  For a fixed tuple
and a fixed \(\varpi\), the equation defining \(\rho\) has at most one
integer solution.  Multiply the counts and use
\(\Theta(N)\ge0.66N/L^2\).  The logarithmic derivative of each term in
(6.3) is negative because its power of \(N\) is negative and the factor
\(L^3\) is dominated by \(N^\delta\) for every displayed
\(\delta>0\) on this half-line.
\end{proof}

\subsection{Exact source reorganization}
\label{sec:source-reorganization}

Let \(\mathcal C\) be a finite cell partition of the two difficult source
domains.  In a cell \(\tau\), let \(X\le p_3p_4<2X\), let
\(\mathcal I_\tau=[M_-,M_+)\) be the interval for
\(m=p_1p_2\rho\), and define its length by
\[
 M=M_+-M_-.
\tag{6.4}
\]
The required grouped coefficients are
\[
\begin{aligned}
 \alpha_{\tau,X}(n)
 &=\sum_{\substack{p_3p_4=n\\(p_3,p_4)\ {\rm in\ the\ }\tau{\rm\ cell}}}
 a_\tau(p_3,p_4),\\
 \beta_{\tau,M}(m)
 &=\sum_{\substack{p_1p_2\rho=m\\(p_1,p_2,\rho)\ {\rm in\ the\ }\tau{\rm\ cell}}}
 b_\tau(p_1,p_2,\rho).
\end{aligned}
\tag{6.5}
\]
The coefficients must satisfy
\[
 \operatorname{supp}\alpha_{\tau,X}\subset[X,2X],\qquad
 \operatorname{supp}\beta_{\tau,M}\subset[M_-,M_+],
\tag{6.6}
\]
\[
 \sum_n|\alpha_{\tau,X}(n)|^2\le12X,\qquad
 \sum_m|\beta_{\tau,M}(m)|^2\le \mathcal E_{\beta,\tau}M.
\tag{6.7}
\]
The cell must make the order condition \(p_2<p_3\), positivity
\(\lceil Z\rceil\le nm\le N-2\), source coprimality, and every range
condition either separable in (6.5) or explicitly charged as
\(E_{\rm src,\tau}\).  Under this stated separability property, finite
distributivity gives the exact identity
\[
\sum_{\substack{(p_1,p_2,p_3,p_4,\rho)\in\tau\\
\lceil Z\rceil\le p_1p_2p_3p_4\rho\le N-2}}
 a_\tau b_\tau
 =
 \sum_{n,m}\alpha_{\tau,X}(n)\beta_{\tau,M}(m)
 \mathbf1_{\{\lceil Z\rceil\le nm\le N-2\}},
\tag{6.8}
\]
with the nonseparable cells appearing in \(E_{\rm src,\tau}\).
Thus (6.8) is an identity only after the partition and the two formulas
(6.5) have been supplied; the words “grouped coefficients” alone do not
establish it.

\subsection{Coprimality and characters}
\label{sec:characters}

Let \(d\) be a supported modulus.  Write
\[
 d=d_Nd_0,\qquad (d_0,N)=1,\qquad
 \text{every prime of }d_N\text{ divides }N.
\tag{6.9}
\]
On the source with \((nm,N)=1\), if \(d_N>1\) then
\(nm\equiv N\pmod {d_N}\) is impossible.  This is a local zero term and
is charged before character expansion.  If \(d_N=1\), then
\[
 \mathbf1_{\{nm\equiv N\pmod d\}}
 =
 \frac1{\varphi(d)}
 \sum_{\chi\pmod d}\chi(nm\overline N),
\tag{6.10}
\]
where all characters are extended by zero off \((\mathbb Z/d\mathbb Z)^*\).
The formula is therefore used only after \((d,N)=1\) and with the
\((nm,d)=1\) restriction displayed.

For each \(\chi\pmod d\), let \(\chi^*\pmod r\) be its primitive inducing
character, \(r\mid d\).  The exact finite grouping is
\[
 \sum_{\chi\pmod d} \chi(nm\overline N)
 =
 \sum_{r\mid d}
 \sum_{\substack{\chi^*\pmod r\\\chi^*\ {\rm primitive}}}
 \sum_{\substack{\chi\pmod d\\\chi\ {\rm induced\ by\ }\chi^*}}
 \chi^*(nm\overline N)
 \mathbf1_{\{(nm,d)=1\}}.
\tag{6.11}
\]
The inner induced-character multiplicity is retained; it is not replaced by
one.  The principal character is separated from the nonprincipal characters,
and an exceptional real character is removed from the nonexceptional sum
and entered into the exceptional ledger.  These are exact bookkeeping
operations, not numerical estimates.

\subsection{The energy lemma with the missing \(+1\)}
\label{sec:energy}

\begin{lemma}[Finite semiprime coefficient energy]
\label{lem:energy}
Let \(\mathcal A\) be a finite set of products \(a=p_1p_2\) of two
distinct primes, let \(w_a\in\mathbb C\), and let \(\mathcal I\) be an
integer interval of length \(M\).  Put
\[
 \beta(m)=\sum_{\substack{a\in\mathcal A\\a\mid m}}w_a,\qquad
 W=\sum_{a\in\mathcal A}|w_a|.
\]
Then
\[
\begin{aligned}
 \sum_{m\in\mathcal I}|\beta(m)|^2
 &\le
 M\sum_{a,b\in\mathcal A}\frac{|w_aw_b|}{[a,b]}+W^2\\
 &=M\left(E_{\rm disj}+E_{\rm share}+E_{\rm diag}\right)+W^2,
\end{aligned}
\tag{6.12}
\]
where
\[
\begin{aligned}
 E_{\rm disj}&=\sum_{(a,b)=1}\frac{|w_aw_b|}{ab},\\
 E_{\rm share}&=\sum_{\substack{(a,b)=p\\p\ {\rm prime}}}
 \frac{|w_aw_b|}{[a,b]},\\
 E_{\rm diag}&=\sum_{a=b}\frac{|w_a|^2}{a}.
\end{aligned}
\tag{6.13}
\]
The first moment obeys
\[
 \sum_{m\in\mathcal I}|\beta(m)|
 \le M\sum_a\frac{|w_a|}{a}+W.
\tag{6.14}
\]
\end{lemma}

\begin{proof}
Expand the square and interchange the finite sums:
\[
 \sum_{m\in\mathcal I}|\beta(m)|^2
 =\sum_{a,b}w_a\overline{w_b}
 \#\{m\in\mathcal I:[a,b]\mid m\}.
\]
An interval of length \(M\) contains at most \(M/[a,b]+1\) multiples
of \([a,b]\), so the triangle inequality gives the first line of (6.12).
For two-prime products, \((a,b)\) is either \(1\), one common prime, or
two common primes.  In the last case the unordered products satisfy
\(a=b\); if ordered pairs are used, the two positional representations
are included in the definition of \(w_a\) and are not counted again.
This proves (6.13).  The same interval count without a second coefficient
gives (6.14).
\end{proof}

\begin{remark}
The frequently written estimate
\(3M(H_{\rm en}^4+4H_{\rm en}^3+2H_{\rm en}^2)\) is not a consequence of
(6.12) until a separate certificate bounds all three sums in (6.13) and
absorbs the \(W^2\) term.  This revision never absorbs \(+1\) by assertion.
\end{remark}

\subsection{The complete conditional bilinear interface}
\label{sec:bilinear}

Define the conductor split parameter and every symbol in the large-conductor
bound as follows:
\[
 2\le q_0<D_*:=Q_N(u_0)D_\circ,\qquad
 K_R=2+\frac{\log(D_*/q_0)}{\log2},
\tag{6.15}
\]
\[
 L_o=\frac{L+\log(3L)}{\pi},\qquad
 L_p=\frac{2L+\log(3L)}{\pi},\qquad
 K_n=2+\frac{(1/3-2\kappa_1)L}{\log2}.
\tag{6.16}
\]
The parameter \(q_0\) is a supplied integer cutoff, not an implicit
conductor.  Its numerical value must be part of a certificate.

\begin{hypothesis}[Complementary-prime switch and bilinear remainder]
\label{hyp:switch}
\(\mathrm{H}_{\rm sw}\) asserts that the cell partition in (6.5) exists,
that (6.8) is exact up to the explicitly listed \(E_{\rm src,\tau}\), and
that after (6.9)--(6.11), Perron separation, and the split
\(r\le q_0\) or \(r>q_0\), the complete signed remainder obeys
\[
 |\mathcal R_{\rm sw}|\le
 \mathcal E_{<q_0}+\mathcal E_{>q_0}
 +\mathcal E_{\rm Perron}+\mathcal E_{\rm local},
\tag{6.17}
\]
where the large-conductor contribution is bounded by the four explicitly
added terms
\[
\begin{aligned}
 \mathcal E_{>q_0}\le&
 1.1\log D_*\sqrt{\mathcal E_\beta}
 \left(B_1+B_2+B_3+B_4\right),\\
 B_1={}&2XM/q_0,\\
 B_2={}&K_R\sqrt{12X}\,M,\\
 B_3={}&K_R\sqrt{12M}\,X,\\
 B_4={}&24D_*\sqrt{XM}.
\end{aligned}
\tag{6.18}
\]
Here \(\mathcal E_\beta\) is the certified right side of (6.12) divided
by \(M\), including its \(W^2/M\) term.  The small-conductor term must
be a stated fixed-target BJS \texttt{PNTPAP1} estimate with its full
variable endpoint; it is not replaced by a prime-in-interval assertion.
\end{hypothesis}

\begin{remark}
The proof obligations behind (6.17) are four separate arrows:
\[
\begin{gathered}
\text{local coprimality}
\longrightarrow\text{character orthogonality}\\
\longrightarrow\text{primitive conductor grouping}
\longrightarrow\text{Cauchy--Schwarz and large sieve}.
\end{gathered}
\]
Equations (6.9)--(6.14) supply the first two arrows and the exact energy
identity.  The inducing multiplicity, the \(+1\) energy term, the
principal/exceptional rows, and the \(q_0,D_*\) ranges still have to be
inserted into a numerical source proof to establish (6.17).  The supplied
manuscript did not contain that proof, so \(\mathrm{H}_{\rm sw}\) remains
\openstatus.  In particular, the corrected (6.18) is a sum, not the
product caused by the missing plus sign in the former display.

This is a genuine reconstruction barrier, not merely an omitted decimal.
In the supplied pre-audit source, \(q_0\) has no numerical definition,
the cells \(\tau\) and the weights \(a_\tau,b_\tau\) are never listed, and
the symbols \(\alpha,\beta\) occur only as \emph{grouped coefficients}.
Thus there is no determinate left side to which the variable-endpoint
\texttt{PNTPAP1} statement can be applied, and no determinate inducing
multiplicity or \(W^2/M\) term to insert into (6.18).  Different choices
of these missing data change both the support and the constant.  Supplying
them would be new mathematical input; they cannot be recovered by
algebraic editing or by interval evaluation of the displayed decimals.
\end{remark}

\subsection{Perron boundaries}
\label{sec:perron}

For \(\sigma>0\), \(T>0\), and \(y>0\), put
\[
 I_T(y)=\frac1{2\pi i}\int_{\sigma-iT}^{\sigma+iT}\frac{y^s}{s}\,ds.
\]
The half-integer identities used to remove equality rows are exact:
\[
 \mathbf1_{\{p_2<p_3\}}
 =\mathbf1_{\{p_3/(p_2+1/2)>1\}},
\tag{6.19}
\]
\[
 \mathbf1_{\{\lceil Z\rceil\le nm\le N-2\}}
 =
 \mathbf1_{\{nm/(\lceil Z\rceil-1/2)>1\}}
 -\mathbf1_{\{nm/(N-3/2)>1\}}.
\tag{6.20}
\]
For \(y\ne1\), direct integration by parts on the two vertical tails gives
\[
 |\,\mathbf1_{\{y>1\}}-I_T(y)\,|
 \le \frac{2y^\sigma}{\pi T|\log y|}.
\tag{6.21}
\]
The minimum logarithmic gaps at the equality rows and the resulting
supported-modulus sum are part of \(\mathcal E_{\rm Perron}\) in (6.17).
No exponent such as \(10^{-10^9}\) is claimed without that sum being
evaluated with outward intervals.

\section{Exceptional conductors and terminal notation}
\label{sec:exceptional}

\subsection{Landau--Page split}
\label{sec:landau}

For \(u>1\), write
\[
 K_\delta(u)=(\log u)^\delta,\qquad Q_1(u)=(\log u)^{10}.
\]
The BJS exceptional alternative gives at most one exceptional primitive
conductor in the relevant range.  If it is below \(K_\delta(x_2(Y))\), the
small-conductor estimate is used.  If it is above that threshold, the
divisible moduli are deleted and the local reciprocal-density factor for
its largest odd prime \(q_1\) is
\[
 \frac{q_1-1}{q_1-2}.
\tag{7.1}
\]
This statement does not imply that an exceptional zero is absent.

\subsection{The endpoint chain calculation}
\label{sec:chain}

Put
\[
 K(c)=(e^c/12)^{1.16},\qquad
 P_{\rm odd}(r)=\prod_{\substack{p\le r\\p>2}}p.
\]
At \(c_0=23.36\), the rational interval certificate in
\cert{final\_interval\_ledger.py} gives
\[
 3.2846409834\cdot10^{10}<K(c_0)
 <3.2846409836\cdot10^{10},
\tag{7.2}
\]
whereas
\[
 P_{\rm odd}(29)=3\,234\,846\,615,\qquad
 P_{\rm odd}(31)=100\,280\,245\,065.
\tag{7.3}
\]
Hence the endpoint cell has \(q_1=31\).  Moreover
\[
 31P_{\rm odd}(23)=3\,457\,939\,485<K(c_0),
\tag{7.4}
\]
so a chain with \(q_2\le23\) cannot reach the endpoint threshold; the
endpoint comparison forces \(q_2\ge29\).

For a fixed primorial cell, the chain coefficient has the form
\[
 \mathcal C(q_1,q_2,\ldots)
 =\frac1{q_1-2}
 \left(1+\frac2{q_2-2}
 +\frac2{(q_2-2)(q_3-2)}+\cdots\right).
\tag{7.5}
\]
The endpoint arithmetic in (7.2)--(7.4) is closed.  The assertion that the
resulting coefficient and all exceptional source charges are no larger on
every later cell is a separate half-line claim and is not inferred from
the endpoint alone.

\begin{hypothesis}[Exceptional and terminal package]
\label{hyp:exceptional}
\(\mathrm{H}_{\rm exc}\) asserts the BJS Landau--Page split, the exceptional
chain bound (7.5) on every cell \(P_k\le K(c)<P_{k+1}\), the
\(\varpi\mid N\) charge (6.2), all \(p=2\), collision, cofactor-one, and
finite Stieltjes endpoint charges, with their required normalization by
\(\Theta(N)\).  It includes both small and large exceptional branches.
\end{hypothesis}

\subsection{Johnston notation}
\label{sec:johnston}

The two previously conflated quantities are now defined separately:
\[
 \frac{T_J^{\rm transfer}}{\Theta(N)}
 \le 4.096650441\cdot10^{-4},
\tag{7.6}
\]
and
\[
 \frac{T_J^{\rm total}}{\Theta(N)}
 =
 \frac{T_J^{\rm transfer}+T_J^{\rm minor}
 +T_J^{\rm collision}+T_J^{\rm endpoint}
 +T_J^{\rm cofactor}}{\Theta(N)}.
\tag{7.7}
\]
The two displayed components
\[
 4.096650441\cdot10^{-4}
 +5.717\cdot10^{-8}
 =4.097222141\cdot10^{-4}
\tag{7.8}
\]
already show that \(4.097\cdot10^{-4}\) is not an upper bound for the
total.  The number \(0.006759\), when used, is a bound for
\(T_J^{\rm total}\), not for (7.6).  The final ledger must charge either
the fully certified total (7.7), or separately certified disjoint
components; it may not charge a component under the name \(T_J\).

\section{The unique ledger}
\label{sec:ledger}

\subsection{Ledger rules}
\label{sec:ledger-rules}

Every ledger row has the fields
\[
\begin{gathered}
(\text{source row},\ \text{Wu coefficient},\ \text{branch factor},\ \text{division by }4),\\
(\text{interval},\ \text{status},\ \text{included charges}).
\end{gathered}
\tag{8.1}
\]
The only accepted normalized master expression is
\[
\begin{aligned}
 \mathcal M_b={}&\frac14\bigl(
 4F_0-F_1-F_2-F_3+F_4+F_5\\
 &\qquad -2F_7-2F_8-F_9-F_{10}
 -U_{13,14}^{(b)}\bigr)\\
 &-E_{\rm bad}/(4\Theta(N))
 -E_{\rm tri}/(4\Theta(N))
 -E_{\rm box}^{(b)}/(4\Theta(N))\\
 &-E_{\rm sw}^{(b)}/(4\Theta(N))
 -T_J^{\rm total,(b)}/\Theta(N)
 -E_{\rm floor}^{(b)}/(4\Theta(N)).
\end{aligned}
\tag{8.2}
\]
Here \(b\) is small or large exceptional; \(U_{13,14}^{(\rm large)}
=(30/29)U_{13,14}^{(\rm small)}\) only when the exceptional package
proves that multiplier.  The old value \(0.616651062059\) is retired and
does not occur in (8.2).  The only retained candidate upper bound is
\[
 U_{13,14}^{(\rm small)}<0.744967,\qquad
 U_{13,14}^{(\rm large)}<\frac{30}{29}\,0.744967.
\tag{8.3}
\]
The factor \(1/4\) in (8.2), the branch factor in (8.3), and the inclusion
or exclusion from \(T_J^{\rm total}\) are recorded on the row itself.

\begingroup
\scriptsize
\setlength{\tabcolsep}{2pt}
\begin{longtable}{p{0.17\textwidth}p{0.09\textwidth}p{0.11\textwidth}p{0.30\textwidth}p{0.25\textwidth}}
\label{tab:ledger-lines}\\
\toprule
row&Wu coefficient&branch factor&normalized candidate charge&accounting rule\\
\midrule
\(U_{13}+U_{14}\)&\(-1/4\)&\(1\), \(30/29\)&
\([0.18624175,0.18624175]\), \([0.1926638792,0.1926638794]\)&
\(U\) bound only; not in \(0.02\)\\
five-box and strip&row weights, already divided by \(4\)&
\(1\), \(30/29\)&
\([0.00346007789,0.00346007791]\),
\([0.00357939092,0.00357939094]\)&
rectangle/geometry charge; not in \(0.02\)\\
\(T_J^{\rm transfer}\)&already normalized&\(1\)&
\([0.00040966504,0.00040966505]\)&
component of (7.7), separately added\\
\(T_J^{\rm minor}\)&already normalized&\(1\)&
\([0.00000005716,0.00000005717]\)&
component of (7.7), separately added\\
remaining finite terminal reserve&already normalized&\(1\)&
\([0.02000000000,0.02000000000]\)&
contains all remaining terminals, excludes the four rows above\\
\(E_{\rm bad},E_{\rm tri},E_{\rm floor}\)&negative&branch-dependent&
\textit{OPEN}&never set to zero\\
\(E_{\rm AP},E_{\rm rect},E_{\rm sw}\)&negative&branch-dependent&
\textit{OPEN}&never set to zero\\
\bottomrule
\end{longtable}
\endgroup
The two rows for \(T_J^{\rm transfer}\) and \(T_J^{\rm minor}\) are not a
claim about \(T_J^{\rm total}\): all other terms in (7.7) must be proved to
fit the separately named \(0.02\) reserve, or \(T_J^{\rm total}\) must be
charged as one larger row.  This convention makes every inclusion decision
visible and prevents double counting.

\subsection{Machine-readable endpoint arithmetic}
\label{sec:ledger-arithmetic}

The exact rational interval script \cert{final\_interval\_ledger.py} uses
the following candidate input intervals, all represented as decimal
rationals:
\[
\begin{array}{c|c}
\text{quantity}&\text{outward interval}\\
\hline
M_{\rm small}^{\rm candidate}
 &[0.86576447960,\ 0.86576447962]\\
M_{\rm large}^{\rm candidate}
 &[0.38981885171,\ 0.38981885173]\\
E_{\rm box}^{\rm small}/\Theta
 &[0.00346007789,\ 0.00346007791]\\
E_{\rm box}^{\rm large}/\Theta
 &[0.00357939092,\ 0.00357939094]\\
U_{13,14}^{\rm small}&[0.7449669,\ 0.7449670]
\end{array}
\tag{8.5}
\]
The intervals in (8.5) are candidate ledger inputs, not proved analytic
outputs.  Exact rational subtraction nevertheless gives
\small
\[
\begin{aligned}
 [0.86576447960,0.86576447962]
 -[0.21011155000,0.21011155020]
 &\subset[0.65565292940,0.65565292962],\\
 [0.38981885171,0.38981885173]
 -[0.21665299240,0.21665299250]
 &\subset[0.17316585921,0.17316585933].
\end{aligned}
\tag{8.6}
\]
\normalsize
The subtraction in (8.6) is closed arithmetic only.  It does not certify
that the two left-hand intervals are the outputs of the Wu, rectangle,
switch, exceptional, and terminal proofs.  The full ledger must reject an
unknown charge rather than silently set it to zero.

\subsection{Critical hypotheses}
\label{sec:critical-hypotheses}

\begin{description}
\item[\(\mathrm{H}_{\rm prod}\)] The finite product lemma (3.2) for every
relevant tag, every required \(u,z,N\), with \(u_0=10^7\).
\item[\(\mathrm{H}_{\rm rect}\)] The global fixed-cutoff rectangle
extension in \cref{hyp:rect}, including \(X\ge N^{0.398}\), all five
boxes, and no per-box \(O(1)\) multiplication.
\item[\(\mathrm{H}_{\rm sw}\)] The exact source regrouping and all terms in
\cref{hyp:switch}, including coprimality, inducing factors, \(+1\) energy,
Perron, and exceptional deletions.
\item[\(\mathrm{H}_{\rm exc}\)] The all-cell exceptional and terminal
package in \cref{hyp:exceptional}.
\item[\(\mathrm{H}_{\rm led}\)] The unique machine-readable ledger proves
that every row in (8.2) has a finite outward upper/lower interval and that
the resulting lower endpoint is at least \(0.1731658\).
\item[\(\mathrm{H}_{\rm half}\)] The interval proof in the next subsection
establishes the same \(0.1731658\) ledger inequality for every
\(c\ge23.36\), including all
primorial, ceiling, branch, and source-domain transitions.
\end{description}

\begin{remark}
\(\mathrm{H}_{\rm prod}\), \(\mathrm{H}_{\rm rect}\), and
\(\mathrm{H}_{\rm sw}\) are the core analysis interfaces.  They are now
precisely stated, but they are not proved by this file.  \(\mathrm{H}_{\rm
exc}\), \(\mathrm{H}_{\rm led}\), and \(\mathrm{H}_{\rm half}\) remain
conditional because they depend on those interfaces and on absent source
certificates.  This is why (8.5) is not promoted to an unconditional
main theorem.

\end{remark}

\section{Half-line uniformity}
\label{sec:halfline}

\subsection{Continuous worst-case envelope}
\label{sec:continuous-envelope}

Define
\[
 \underline\theta(c)=
 \frac12-\frac{\log Q_{\max}(10^7)+20c}{e^c}.
\tag{9.1}
\]
Since \(Q_N(10^7)\le Q_{\max}(10^7)\), any level expression with
\(Q_N(10^7)\) is bounded below by this fixed envelope.  Its derivative is
\[
 \underline\theta'(c)=e^{-c}
 \bigl(\log Q_{\max}(10^7)+20c-20\bigr)>0
\quad(c\ge23.36).
\tag{9.2}
\]
This derivative treats \(Q_{\max}\) as a fixed number; it does not
differentiate the discontinuous \(Q_N\).

For the boxes,
\[
 M_i(c)\le
 \frac{(b_i-a_i)e^c}{\log(1+\eps_0)}+1.
\tag{9.3}
\]
Any claimed monotonicity of a normalized error must be proved for the full
product containing the right side of (9.3), not for a single factor.
Likewise the energy symbol \(H_{\rm en}\) is not replaced by the level
function \(H_{\rm lev}(c)\).

\subsection{Discrete-cell protocol}
\label{sec:cell-protocol}

The primorial cells are
\[
 P_k\le K(c)<P_{k+1},\qquad
 c\in\left[
 \log12+\frac{\log P_k}{1.16},\
 \log12+\frac{\log P_{k+1}}{1.16}
 \right).
\tag{9.4}
\]
On each such cell the chain indices are fixed.  The ceiling transitions
of (9.3), the source-domain endpoints, and the two exceptional alternatives
are added to the partition.  On every finite interval \(I\) the required
certificate is
\[
 \inf_{c\in I}
 \left(\mathcal M_b(c)-E_{\rm all}^{(b)}(c)\right)>0.17,
\tag{9.5}
\]
with all integrals and logarithms evaluated by outward intervals.  At each
transition both one-sided limits are checked.  On the final half-infinite
interval an analytic inequality must dominate every remaining polynomial
factor by its negative exponential factor.

\begin{hypothesis}[Half-line interval cover]
\label{hyp:half}
\(\mathrm{H}_{\rm half}\) is the existence of a finite or finitely described
cover satisfying (9.4)--(9.5), with a supplied interval record for every
cell.  It must include the box-count envelope, \(Q_{\max}\), the exceptional
chain, all positive and negative signed terms, and every rounding boundary.
\end{hypothesis}

\begin{remark}
The old endpoint differentiation did not establish (9.5): it treated
\(Q_N\) as constant, used an energy quantity under the name \(H\), and
ignored box-count and primorial jumps.  The endpoint \(c_0\) is therefore
not asserted to be the worst point in this revision.  This explicitly
closes the logical route from “two endpoint values” to “all \(c\ge c_0\)”
only by making the missing proof a named hypothesis.
\end{remark}

\section{Proof-status matrix}
\label{sec:status-matrix}

\begingroup
\scriptsize
\setlength{\tabcolsep}{2pt}
\begin{longtable}{p{0.11\textwidth}p{0.45\textwidth}p{0.25\textwidth}p{0.10\textwidth}}
\toprule
node&obligation&where handled&status\\
\midrule
W&Wu signed identity and exact row signs&\cref{sec:wu-input}&quoted\\
W-bad&positive high-multiplicity rows removed before compensation&
\cref{lem:wu-corrected}&\closedstatus\\
R&full-product coprimality switch and \(\varpi\mid N\) terminal&
\cref{sec:rough-switch,lem:varpi-N}&\closedstatus\\
P&\(u_0=10^7\) finite product and Abel tail&
\cref{sec:products,hyp:prod}&\openstatus\\
T&tag-dependent \(X_A\), \(V_\alpha\), and product quantifiers&
\cref{lem:tagged}&\closedstatus\\
B&five-box positivity, endpoints, and exact count&
\cref{sec:box-definition,lem:boxes}&\closedstatus\\
Rect&\(X\ge N^{0.398}\) global rectangle estimate&
\cref{hyp:rect}&\openstatus\\
E&semiprime energy with diagonal/shared/disjoint cases and \(+1\)&
\cref{lem:energy}&\closedstatus\\
Sw&source regrouping, character, conductor, Perron, large sieve&
\cref{hyp:switch}&\openstatus\\
J&transfer versus total Johnston terminal&\cref{sec:johnston}&\closedstatus\textsuperscript{*}\\
L&single ledger, all charges and branch factors&
\cref{sec:ledger}&\openstatus\\
H&all-\(c\) interval cover and transition checks&
\cref{sec:halfline}&\openstatus\\
\bottomrule
\end{longtable}
\endgroup

\noindent\textsuperscript{*}\(\closedstatus\) means the notation and
accounting rule are corrected; the numerical total remains an input to
\(\mathrm{H}_{\rm exc}\) or \(\mathrm{H}_{\rm led}\), not an unconditional
claim.

\section{Reproducible interval protocol}
\label{sec:protocol}

\subsection{Rational outward arithmetic}
\label{sec:interval}

Every displayed finite decimal in the candidate ledger is interpreted as a
rational interval \([a,b]\) with decimal endpoints.  For positive intervals
\[
 [a,b][c,d]=[\min(ac,ad,bc,bd),\max(ac,ad,bc,bd)],
\]
and subtraction is
\[
 [a,b]-[c,d]=[a-d,b-c].
\]
The script does not use binary floating point for the ledger assertions.
For logarithms it records a rational input enclosure; when a new logarithm
is generated, it may use
\[
 \log x=2\sum_{j=0}^{m}\frac{u^{2j+1}}{2j+1}+R_m,\qquad
 u=\frac{x-1}{x+1},
\]
with
\[
 0<R_m<
 \frac{2u^{2m+3}}{(2m+3)(1-u^2)}
\quad(0<u<1).
\tag{11.1}
\]
Products, divisions, and the ceiling tests in (5.11) are then exact
rational operations.  No output labelled \textit{numerical check} is
accepted as a proof of a missing analytic inequality.

\subsection{Certificate manifest}
\label{sec:manifest}

The supplied files for this revision are:
\[
\begin{array}{ll}
\cert{explicit\_chen\_c2336\_final.tex}&
\text{this complete conditional manuscript},\\
\cert{final\_ledger.json}&
\text{one machine-readable candidate ledger},\\
\cert{final\_interval\_ledger.py}&
\text{exact rational arithmetic and box-count checks},\\
\cert{final\_audit.md}&
\text{item-by-item audit disposition}.
\end{array}
\tag{11.2}
\]
The script verifies (5.11), (5.12), (7.2)--(7.4), and the arithmetic
subtractions in (8.6).  It intentionally fails closed when an analytic
interface is marked \openstatus; it never reports an unconditional
positive Chen margin.

\section{Final statement}
\label{sec:final}

The corrected file has removed the false \(N\)-independent rough-number
claim, the undefined \(\eps\), \(q_0\), \(D_*\), \(M\), \(H\), \(B_\circ\),
and \(F_+\) conventions, the missing plus signs, the wrong box count, the
negative-source definition, and the conflated Johnston bounds.  It has also
made the full-product switch, floor term, local non-coprime branches,
tagged local products, energy \(+1\), and half-integer boundaries explicit.

It has not supplied proofs of \(\mathrm{H}_{\rm prod}\),
\(\mathrm{H}_{\rm rect}\), or \(\mathrm{H}_{\rm sw}\), nor the dependent
all-cell ledger and half-line cover.  Therefore the strongest statement
that can be made without inventing an external theorem or a missing
certificate is exactly \cref{thm:conditional-conclusion} and
\cref{prop:status}, namely:
\[
\boxed{\text{\rm Class 3, conditional reduction theorem only.}}
\]

\begin{thebibliography}{9}
\bibitem{BJS}
M. Bordignon, D. R. Johnston and V. Starichkova,
\emph{An explicit version of Chen's theorem and the linear sieve},
\emph{International Journal of Number Theory} 21 (2025), 2497--2572;
\href{https://arxiv.org/abs/2207.09452}{arXiv:2207.09452v6}.
\bibitem{WuI}
J. Wu, \emph{Chen's double sieve, Goldbach's conjecture and the twin prime
problem}, \emph{Acta Arithmetica} 114 (2004), 215--273;
\href{https://arxiv.org/abs/0705.1652}{arXiv:0705.1652}.
\end{thebibliography}

\end{document}

